
\documentclass{article}
\usepackage{colortbl}
\usepackage{makecell}
\usepackage{multirow}
\usepackage{supertabular}

\begin{document}

\newcounter{utterance}

\centering \large Interaction Transcript for game `cladder', experiment `full\_v1.5\_default', episode 5313 with LLP: Large Language Problems\_\_qwen3.5{-}9b\_\_llp{-}final.
\vspace{24pt}

{ \footnotesize  \setcounter{utterance}{1}
\setlength{\tabcolsep}{0pt}
\begin{supertabular}{c@{$\;$}|p{.15\linewidth}@{}p{.15\linewidth}p{.15\linewidth}p{.15\linewidth}p{.15\linewidth}p{.15\linewidth}}
    \# & \multicolumn{2}{c}{Player} && \multicolumn{2}{c}{Game Master} \\
    \hline

    \theutterance \stepcounter{utterance}  
    & & & \multicolumn{4}{p{0.6\linewidth}}{
        \cellcolor[rgb]{0.9,0.9,0.9}{
            \makecell[{{p{\linewidth}}}]{
                \texttt{\tiny{[P1$\langle$GM]}}
                \texttt{You are an expert in causal inference. The following question is not a typical commonsense query, but rather a meticulously designed question created by a professor specializing in causal inference, intended to assess the students' mastery of the course content.} \\
\texttt{Imagine a self{-}contained, hypothetical world with only the following conditions, and without any unmentioned factors or causal relationships: Liking spicy food has a direct effect on occupation and salary. Skill has a direct effect on occupation and salary. Occupation has a direct effect on salary. Skill is unobserved. The overall probability of liking spicy food is 2\%. The probability of not liking spicy food and high salary is 45\%. The probability of liking spicy food and high salary is 1\%. Is the chance of high salary smaller when observing liking spicy food?} \\
\texttt{Start your answer with 'Yes' or 'No', followed by additional reasoning or evidence to support your explanation.} \\
            }
        }
    }
     \\ \\

    \theutterance \stepcounter{utterance}  
    & \multicolumn{4}{p{0.6\linewidth}}{
        \cellcolor[rgb]{0.561,0.737,0.561}{
            \makecell[{{p{\linewidth}}}]{
                \texttt{\tiny{[P1$\rangle$GM]}}
                \texttt{No.} \\
\\ 
\texttt{**Reasoning:**} \\
\\ 
\texttt{To determine whether the chance of high salary is smaller when observing liking spicy food, we need to compare the conditional probability of high salary given liking spicy food, \$P(\textbackslash{}text\{High Salary\} \textbackslash{}mid \textbackslash{}text\{Liking Spicy\})\$, with the overall probability of high salary, \$P(\textbackslash{}text\{High Salary\})\$.} \\
\\ 
\texttt{We are given:} \\
\texttt{{-} \$P(\textbackslash{}text\{Liking Spicy\}) = 0.02\$} \\
\texttt{{-} \$P(\textbackslash{}text\{Not Liking Spicy\} \textbackslash{}cap \textbackslash{}text\{High Salary\}) = 0.45\$} \\
\texttt{{-} \$P(\textbackslash{}text\{Liking Spicy\} \textbackslash{}cap \textbackslash{}text\{High Salary\}) = 0.01\$} \\
\\ 
\texttt{First, compute the overall probability of high salary:} \\
\texttt{\$\$} \\
\texttt{P(\textbackslash{}text\{High Salary\}) = P(\textbackslash{}text\{Liking Spicy\} \textbackslash{}cap \textbackslash{}text\{High Salary\}) + P(\textbackslash{}text\{Not Liking Spicy\} \textbackslash{}cap \textbackslash{}text\{High Salary\}) = 0.01 + 0.45 = 0.46} \\
\texttt{\$\$} \\
\\ 
\texttt{Now compute the conditional probability:} \\
\texttt{\$\$} \\
\texttt{P(\textbackslash{}text\{High Salary\} \textbackslash{}mid \textbackslash{}text\{Liking Spicy\}) = \textbackslash{}frac\{P(\textbackslash{}text\{Liking Spicy\} \textbackslash{}cap \textbackslash{}text\{High Salary\})\}\{P(\textbackslash{}text\{Liking Spicy\})\} = \textbackslash{}frac\{0.01\}\{0.02\} = 0.5} \\
\texttt{\$\$} \\
\\ 
\texttt{Compare:} \\
\texttt{{-} \$P(\textbackslash{}text\{High Salary\} \textbackslash{}mid \textbackslash{}text\{Liking Spicy\}) = 0.5\$} \\
\texttt{{-} \$P(\textbackslash{}text\{High Salary\}) = 0.46\$} \\
\\ 
\texttt{Since \$0.5 > 0.46\$, the chance of high salary is **larger**, not smaller, when observing liking spicy food.} \\
\\ 
\texttt{Therefore, the answer is **No**.} \\
            }
        }
    }
    & & \\ \\

    \theutterance \stepcounter{utterance}  
    & & & \multicolumn{2}{p{0.3\linewidth}}{
        \cellcolor[rgb]{0.9,0.9,0.9}{
            \makecell[{{p{\linewidth}}}]{
                \texttt{\tiny{[GM$|$GM]}}
                \texttt{No} \\
            }
        }
    }
    & & \\ \\

    \theutterance \stepcounter{utterance}  
    & & & \multicolumn{2}{p{0.3\linewidth}}{
        \cellcolor[rgb]{0.9,0.9,0.9}{
            \makecell[{{p{\linewidth}}}]{
                \texttt{\tiny{[GM$|$GM]}}
                \texttt{yes} \\
            }
        }
    }
    & & \\ \\

    \theutterance \stepcounter{utterance}  
    & & & \multicolumn{2}{p{0.3\linewidth}}{
        \cellcolor[rgb]{0.9,0.9,0.9}{
            \makecell[{{p{\linewidth}}}]{
                \texttt{\tiny{[GM$|$GM]}}
                \texttt{game\_result = LOSE} \\
            }
        }
    }
    & & \\ \\

\end{supertabular}
}

\end{document}
